% ============================================================
\section{EvoSuite}
\label{sec:evosuite}
% ============================================================

EvoSuite genera una test suite tramite algoritmo genetico che opera sui casi di test stessi (non
sul codice sorgente): la fitness combina la copertura del criterio scelto con la compattezza della
suite.

\subsection{BrokerImpl}

Versione 1.2.0. Su Maven Central la versione più recente è la 1.0.6, incompatibile con il bytecode Java 11
con cui è compilato il progetto; anche il plugin Maven
ufficiale, fermo alla stessa 1.0.6, è incompatibile: da qui la scelta del jar standalone. Criterio di copertura lasciato al default
dello strumento, che ottimizza contemporaneamente più criteri
(\texttt{LINE;BRANCH;EXCEPTION;WEAKMUTATION;OUTPUT;METHOD;METHODNOEXCEPTION;CBRANCH})
(\hyperref[cmd:evosuite-brokerimpl-gen]{Appendice}).

Un primo tentativo con budget di 60 secondi aveva dato 43,7\% di copertura aggregata; portando il
budget a 300 secondi la copertura sale al 48,7\%, con 456 metodi
\texttt{@Test} generati in 303 secondi. Log completo in
\texttt{report/data/evosuite-brokerimpl/evosuite-log.txt}, statistiche in
\texttt{report/data/evosuite-brokerimpl/evosuite-report/statistics.csv}.

A differenza di Randoop e Test LLM, la suite non è stata integrata nel test tree Maven. I test
generati richiedono a runtime, non solo in generazione, le classi di \texttt{org.evosuite.runtime},
pacchettizzate nella libreria
\texttt{evosuite-standalone-runtime}. Su Maven Central questa libreria si ferma anch'essa alla
versione 1.0.6, incompatibile con i test generati con la 1.2.0. La suite resta quindi documentata,
compilata ed eseguita a parte con la stessa libreria scaricata manualmente
(\hyperref[cmd:evosuite-brokerimpl-offline]{Appendice}).

452 test su 456 passano. L'output di \texttt{JUnitCore} su questa esecuzione
(\hyperref[log:evosuite-offline-junitcore]{Appendice}) nomina i 4 fallimenti e ne mostra lo stack
trace: \texttt{test097} fallisce per una risorsa di servizio (\texttt{ProductDerivation}) mancante
nella classpath minimale usata per questa sola verifica offline; \texttt{test098} fallisce a
cascata per lo stesso motivo. \texttt{test000} e \texttt{test137} falliscono entrambi con la stessa
eccezione (\texttt{GeneralException: Mock call to processArgument...}).

\subsection{LoginDialog}
% non affrontato in questo giro, per decisione presa in Sezione~\ref{sec:randoop}
